\section{Interior buffers and boundary orientation}
\label{sec:local}

The exact semigroup has $a_{\rm ex}=\e^{-x}$ and $|c_{\rm ex}|^2=\e^{-x-2u}$, hence
\begin{equation}
 F_{{\rm ex},\theta}=\e^{-x}(1-\e^{-2u})
 +\theta(1-\theta)(1-\e^{-x})^2.
 \label{eq:exact_buffer}
\end{equation}
Positive dephasing contributes at first order, while bidirectional relaxation contributes at second order. The one-way, zero-dephasing endpoints are singular because both buffers vanish.

\subsection{Local interior laws}

Let $R(z)=1+z+O(z^2)$ be consistent. As $(x,u,\nu)\to0$,
\begin{equation}
 \Fmargin(x,u,\nu)
 =2u+O(x^2+xu+u^2+\nu^2).
 \label{eq:dephasing_buffer}
\end{equation}
Along a fixed generator ray
\begin{equation}
 u=\kappa x,
 \qquad \nu=\varpi x,
 \qquad \kappa=\gamma_\phi/\Gamma,
 \qquad \varpi=\omega/\Gamma,
 \label{eq:generator_ray}
\end{equation}
any $\kappa>0$ supplies a nonzero local CPTP interval.

When $u=0$, write $R(z)=1+z+r_2z^2+O(z^3)$. Expansion of Eq.~\eqref{eq:full_margin} gives
\begin{equation}
 \Fmargin(x,0,\varpi x)
 =C_{R,\theta}(\varpi)x^2+O(x^3),
 \label{eq:excitation_expansion}
\end{equation}
where
\begin{equation}
 C_{R,\theta}(\varpi)
 =\theta(1-\theta)+\frac{r_2}{2}-\frac14+(2r_2-1)\varpi^2.
 \label{eq:excitation_coeff}
\end{equation}
Every method of order at least two has $r_2=1/2$, so
\begin{equation}
 \Fmargin(x,0,\varpi x)
 =\theta(1-\theta)x^2+O(x^3)>0
 \label{eq:excitation_buffer}
\end{equation}
for $0<\theta<1$ and sufficiently small $x$. Dephasing and bidirectional relaxation move the exact trajectory into the relative interior of the relevant channel face. The signed boundary theory below therefore applies precisely at $\theta\in\{0,1\}$ and $u=0$.

On the pure-dephasing/precession boundary $x=0$, the criterion is
\begin{equation}
 |R(-u+\ii\nu)|\le1.
 \label{eq:pure_dephasing}
\end{equation}
The purely unitary slice is governed by imaginary-axis absolute stability, $|R(\ii\nu)|\le1$ \cite{HairerNorsettWanner1993,IserlesNorsett1991OrderStars}.

For later controller estimates, specialize to RK4 and a rotation-dominated joint limit in which $\kappa\to0$ or $\theta(1-\theta)\to0$, $|\varpi|\to\infty$, and the resulting attached endpoint satisfies $x\to0$ and $|\varpi|x\ll1$. The outward boundary defect is
\begin{equation}
 M_4(x;\varpi)=-\frac{\varpi^4x^5}{24}+o(\varpi^4x^5).
\end{equation}
Balancing it against the local buffers gives
\begin{equation}
 x_{\rm att}^{(\phi)}\sim\frac{(48\kappa)^{1/4}}{|\varpi|},
 \qquad
 x_{\rm att}^{(\theta)}\sim
 \left[\frac{24\theta(1-\theta)}{\varpi^4}\right]^{1/3}.
 \label{eq:attached_scaling}
\end{equation}
These are joint-limit estimates, not finite-frequency identities.

\subsection{Signed first-defect law}

Fix the one-way, zero-dephasing boundary and set $\alpha=1/2-\ii\varpi$. Suppose
\begin{equation}
 R(z)-\e^z=\eta_m z^m+O(z^{m+1}),
 \qquad m\ge2,
 \qquad \eta_m\ne0.
 \label{eq:first_defect}
\end{equation}
Then $|\e^{-\alpha x}|^2=\e^{-x}$ cancels all lower-order semigroup terms.

\begin{theorem}[Frequency-dependent boundary law]
\label{thm:defect}
Under Eq.~\eqref{eq:first_defect},
\begin{equation}
 \Mmargin(x)
 =(-1)^m\eta_mG_m(\varpi)x^m+O(x^{m+1}),
 \label{eq:defect_law}
\end{equation}
where
\begin{equation}
 G_m(\varpi)=1-2\operatorname{Re}(1/2-\ii\varpi)^m.
 \label{eq:G_definition}
\end{equation}
For sufficiently small $x>0$, the sign of $(-1)^m\eta_mG_m$ determines whether the RK map enters or exits the channel set, except when $G_m=0$.
\end{theorem}

Define
\begin{equation}
 q=1+4\varpi^2,
 \qquad
 \sigma_m=2^m(\alpha^m+\bar\alpha^m).
 \label{eq:sigma_definition}
\end{equation}
Because $\alpha$ and $\bar\alpha$ solve $z^2-z+q/4=0$,
\begin{equation}
 \sigma_0=\sigma_1=2,
 \qquad
 \sigma_m=2\sigma_{m-1}-q\sigma_{m-2},
 \qquad
 G_m=1-\frac{\sigma_m}{2^m}.
 \label{eq:G_recurrence}
\end{equation}
For the methods used below,
\begin{align}
 G_4&=-\frac18q(q-8),
 &G_5&=-\frac5{16}q(q-4),\nonumber\\
 G_6&=\frac1{32}q(q^2-18q+48),
 &G_7&=\frac7{64}q(q-4)^2.
 \label{eq:G_factors}
\end{align}
At $q=4$, or $\varpi=\sqrt3/2$, $\alpha=\e^{-\ii\pi/3}$ and
\begin{equation}
 G_m(\sqrt3/2)=1-2\cos(m\pi/3).
 \label{eq:unit_circle_rule}
\end{equation}
Thus the unit-circle cancellation occurs when $m\equiv\pm1\pmod6$. A simple zero reverses orientation; an even-multiplicity zero produces tangency.

\begin{proposition}[Degenerate frequency]
\label{prop:degenerate}
If $R(z)-\e^z=\eta_mz^m+\eta_{m+1}z^{m+1}+O(z^{m+2})$ and $G_m(\varpi)=0$, then
\begin{align}
 M_{R,\varpi}(x)
 ={}&\Big[(-1)^{m+1}\eta_{m+1}G_{m+1}(\varpi)\nonumber\\
 &+2(-1)^m\eta_m\operatorname{Re}(\bar\alpha\alpha^m)\Big]x^{m+1}
 +O(x^{m+2}).
 \label{eq:degenerate_coefficient}
\end{align}
\end{proposition}

For SSPRK(3,3) and classical RK4 the cancellation cases factor exactly:
\begin{align}
 M_{{\rm SSP3},\sqrt7/2}(x)&=\frac{x^5(3-2x)}{9},\nonumber\\
 M_{{\rm RK4},\sqrt3/2}(x)&=-\frac{x^6(x-2)^2}{576}.
 \label{eq:threshold_factors}
\end{align}
The principal Dormand--Prince fifth-order formula \cite{DormandPrince1980} has two reversal frequencies,
\begin{equation}
 \varpi=\sqrt{2-\frac{\sqrt{33}}4},
 \qquad
 \varpi=\sqrt{2+\frac{\sqrt{33}}4},
 \label{eq:dp5_thresholds}
\end{equation}
so its small-step orientation is inward, outward, then inward. Formal RK order alone cannot predict complete positivity.
